\subsection{Final Factorial Evaluation}

The final evaluation comprised 16 held-out CT volumes representing 14 independent
source-subject clusters. Across the four outer folds, six sparse-supervision
conditions were evaluated using three independently trained seeds, yielding 288
frozen probability maps. The final evaluator and all operating-point rules were
prospectively frozen before the dense outer-CV references were accessed.

\subsubsection{Primary Low-False-Positive Endpoint}

The prospectively defined primary endpoint was macro-patient lesion recall at
no more than one false positive per patient (R@1). All six factorial conditions
obtained an R@1 of 0.000. The 10,000-replicate paired source-subject cluster
bootstrap yielded a 95\% confidence interval of [0.000, 0.000] for every
condition. The same floor effect was observed at the additional predeclared
budgets of 0.5, 2, and 4 false positives per patient.

The selected operating threshold for every condition and seed was the predefined
zero-prediction fallback threshold. Consequently, none of the tested sparse
annotation configurations achieved non-zero lesion recall while satisfying the
prospectively specified low-false-positive operating constraints.

\subsubsection{Primary Factorial Interactions}

The two predeclared coverage-by-geometry interactions were therefore also
identically zero. The dispersed-geometry interaction estimate was 0.000 and the
fragmented-geometry interaction estimate was 0.000. Exact two-sided
source-subject cluster sign-flip tests produced raw $p$-values of 1.000 for both
contrasts. After Holm correction across the two primary interaction tests, both
adjusted $p$-values remained 1.000. Thus, the primary endpoint provided no
evidence of a coverage-by-geometry interaction under the frozen low-FP criterion.

\subsubsection{Descriptive Secondary Segmentation Behaviour}

Although the primary lesion-level operating point collapsed to zero, the
fixed-threshold secondary analysis showed that the probability maps retained
lesion-related spatial signal. At a prospectively fixed probability threshold of
0.50, any-overlap lesion recall ranged from approximately 0.486 to 0.664 across
conditions. In contrast, recall requiring an IoU of at least 0.25 ranged only
from approximately 0.026 to 0.212. Mean Dice remained between approximately
0.034 and 0.059, and mean IoU remained between approximately 0.018 and 0.031.

The divergence between lesion contact and geometrically valid lesion recovery
was accompanied by substantial prediction fragmentation and false-positive
burden. Mean prediction-component counts ranged from approximately 2,677 to
8,051 components per evaluated volume, while mean false-positive volume exceeded
11,600~mL in every condition. These secondary endpoints are reported as
descriptive mechanistic analyses and were not substituted for the prospectively
defined primary endpoint.

Among the descriptive secondary measures, C50--COH exhibited the highest
any-overlap lesion recall and the highest IoU$\geq$0.25 lesion recall. C100--FRG
showed the highest mean Dice and IoU and the lowest mean false-positive volume,
whereas C100--COH showed the lowest mean fragmentation ratio. These observations
describe the behaviour of the tested conditions but are not interpreted as
statistically established superiority because no confirmatory secondary
hypothesis-testing procedure was prospectively defined.
